\documentclass[12pt,a4paper]{article}
\usepackage[utf8]{inputenc}
\usepackage{amsmath, amssymb, amsfonts}
\usepackage{geometry}
\geometry{margin=2.5cm}
\usepackage{hyperref}
\usepackage{pgfplots}
\usepackage{fancyhdr}
\usepackage{listing}
\usepackage{listings}
\usepackage{titlesec}
\titleformat{\section}
  {\normalfont\Large\bfseries\centering}  % Format: font size, bold, centered
  {\thesection}{1em}{}
\pagestyle{fancy}
\pgfplotsset{compat=1.18}
\fancyhead{}
\fancyfoot{}
\title{Architetture e Sistemi Operativi}
\author{}
\date{}
\fancyhead[L]{Gabriele Perna}          % Left side
\fancyhead[C]{Architetture e Sistemi Operativi}     % Center
\fancyhead[R]{2026/2027}             % Right side
\fancyfoot[C]{\thepage}           % Page number in the center of footer
\let\oldsection\section
\renewcommand{\section}{\clearpage\oldsection}

\begin{document}

\maketitle
\tableofcontents

\section{Introduzione}
In questo corso andremo a:
\begin{itemize}
    \item Capire come funziona un elaboratore.
    \item Capire quali tipi di istruzioni è in grado di eseguire un elaboratore nativamente.
    \item Come lo si può dotare di un sistema operativo.
\end{itemize}
Il corso utilizza Windows e Linux come esempi di sistemi operativi.
\paragraph{Codice Teams:} 6uk5ntg
\paragraph{Email Docente:} \href{marco.delutto@unipi.}{marco.delutto@unipi.}

\section{Processore, memoria e istruzioni}
Il \textbf{processore} è in grado di fare operazioni piccole.\\
Il processore contiene i \textbf{registri}, memorie piccolissime, e il \textbf{Program Counter}, che conserva lo stato di esecuzione memorizzando l'indirizzo della prossima istruzione da eseguire..\\
Il processore è unito ad una \textbf{memoria}, costituita da \textbf{celle}, è vista come un vettore:
\[
M[i] = k
\]
Ogni cella ha un valore e la sua capacità è 1 parola (32 o 64 bit).\\
Una volta che il processore viene acceso, esso esegue:
\begin{lstlisting}
    while(true):
        instr = getInstructionFrom(M)
        decoded = decode(instr)
        execute(decoded)
        determineResults()
        \dots
\end{lstlisting}
Il processore quindi:
\begin{itemize}
    \item Estrae istruzioni dalla memoria
    \item Le decodifica
    \item Le esegue
    \item Calcola i risultati
\end{itemize}

\subsection{Tipi di istruzioni}
Le istruzioni possono essere:
\begin{itemize}
    \item Operative
    \begin{itemize}
        \item ADD
        \item SUB
        \item MUL
        \item DIV
        \item AND
        \item OR
        \item NOT
        \item SHIFT
        \item ROT
        \item \dots
    \end{itemize}
    \item Di Memoria
    \begin{itemize}
        \item LOAD
        \begin{itemize}
            \item Restituisce un valore da un indirizzo:\\
            \begin{lstlisting}
                LOAD val,addr
            \end{lstlisting}
        \end{itemize}
        \item STORE
        \begin{itemize}
            \item Salva un valore in un indirizzo:\\
            \begin{lstlisting}
                STORE val,addr
            \end{lstlisting}
        \end{itemize}
    \end{itemize}
    \item di Salto
    \begin{itemize}
        \item "Saltano" pezzi di codice, volendo anche secondo istruzioni.\\
        \begin{lstlisting}
            inizio:
            \dots
            jmp fine:
            -skipped code-
            fine:
        \end{lstlisting}
    \end{itemize}
\end{itemize}

\end{document}